% 共通の組版設定。本文と証明は各論文の節ファイルに一度だけ置く。
\documentclass[11pt,a4paper]{ltjsarticle}
\usepackage[margin=24mm]{geometry}
\usepackage{amsmath,amssymb,amsthm,mathtools,booktabs,longtable}
\usepackage{mathrsfs}
\usepackage[haranoaji]{luatexja-preset}
\usepackage[unicode,hidelinks]{hyperref}
\setlength{\emergencystretch}{3em}
\hypersetup{bookmarksnumbered=true}
\pdfstringdefDisableCommands{\def\qgr{qgr}\def\HH{HH}\def\Perf{Perf}}
\providecommand{\PublicationVersion}{publication-v32}
\newtheorem{theorem}{定理}[section]
\newtheorem{proposition}[theorem]{命題}
\newtheorem{lemma}[theorem]{補題}
\newtheorem{corollary}[theorem]{系}
\theoremstyle{definition}\newtheorem{definition}[theorem]{定義}
\newtheorem{assumption}[theorem]{入力・仮定}
\theoremstyle{remark}\newtheorem{remark}[theorem]{注意}
\renewcommand{\proofname}{証明}
\newcommand{\HH}{\operatorname{HH}}
\newcommand{\qgr}{\operatorname{qgr}}
\newcommand{\coh}{\operatorname{coh}}
\newcommand{\Perf}{\operatorname{Perf}}
\newcommand{\Ext}{\operatorname{Ext}}
\newcommand{\Hom}{\operatorname{Hom}}
\newcommand{\Spec}{\operatorname{Spec}}
\newcommand{\Spf}{\operatorname{Spf}}
\newcommand{\GL}{\operatorname{GL}}
\newcommand{\gldim}{\operatorname{gldim}}
\newcommand{\OO}{\mathcal O}
\renewcommand{\AA}{\mathbb A}
\newcommand{\PP}{\mathbb P}

% 各原稿の引用だけを共通文献データから組版する。
\makeatletter
\newcommand{\UsePublicationReferences}[1]{%
  \@for\pub@reference:=#1\do{%
    \expandafter\def\csname pubref@\pub@reference\endcsname{}%
  }%
}
\makeatother
